\begin{tabular}{p{0.20\textwidth}p{0.28\textwidth}p{0.24\textwidth}p{0.22\textwidth}}
\toprule
Метрика & Как читать & Что получилось & Интерпретация \\
\midrule
Относительное улучшение & 5--10\% — небольшое; 10--20\% — содержательное; 20--30\% — сильное; 50\%+ — очень большое. & Для обучения по состоянию против фиксированного правила среднее улучшение = 61.4\%; по средним сценариев = 49.9--67.2\%. & Очень большой выигрыш относительно фиксированного правила. \\
Доля побед & 0.5 — преимущества нет; 0.7--0.8 — хорошо; 0.8--0.9 — сильно; 1.0 — полное доминирование на выборке. & Для обучения по состоянию против фиксированного правила доля побед = 1.0 во всех 4 сценариях. & На проверочных траекториях фиксированное правило не выигрывает ни в одном сценарии. \\
Вероятность ухудшения & Чем ближе к нулю, тем безопаснее улучшение. & Для обучения по состоянию против фиксированного правила вероятность ухудшения = 0.0. & На проверенных траекториях ухудшения не наблюдается. \\
Стандартизованный эффект & 0.2 — маленький; 0.5 — средний; 0.8+ — большой. & Парный стандартизованный эффект для обучения по состоянию против фиксированного правила = 1.34. & Большой эффект относительно собственного разброса. \\
Хвостовой риск & Важно смотреть, не исчезает ли выигрыш в худших эпизодах. & 90-й процентиль разности потерь = -9.039e-04; средняя разность в худших 10\% случаев = -4.155e-04. & Даже в верхнем хвосте разность остается ниже нуля, значит эффект не держится на удачных выбросах. \\
Покомпонентный механизм & Показывает, за счет чего именно уменьшается функция потерь. & Среднее снижение компоненты инфляции = 64.8\%, выпуска = 65.1\%, сглаживания ставки = 55.2\%. & Выигрыш идет через лучшую стабилизацию инфляции и выпуска, а не только через механическое сглаживание ставки. \\
Плавность ставки & Если качество лучше при меньшей волатильности, политика не покупает выигрыш чрезмерной дерганостью. & Среднее снижение волатильности ставки для обучения по состоянию против фиксированного правила = 43.5\%. & Обучаемое правило в среднем и лучше, и плавнее фиксированного. \\
Переносимость на новые среды & Коэффициент переноса, близкий к 1, означает, что выигрыш почти сохраняется вне базовой среды. & Относительно фиксированного правила переносимый выигрыш = 51.9--61.6\%; коэффициент переноса = 0.84--1.00. & Эффект хорошо переносится относительно фиксированного правила. \\
Граница применимости & Важно проверить, не является ли выигрыш универсальным по отношению ко всем альтернативам. & Прямая политика по наблюдениям выигрывает у обучения по состоянию только в 27\% проверочных траекторий; при переносе на новые среды обучаемое правило хуже перенастроенного простого правила на 326.4--477.0\%. & Обучаемая политика полезна не универсально, а прежде всего как альтернатива жесткому фиксированному правилу и ошибочной фильтрации. \\
\bottomrule
\end{tabular}